\documentclass{article}
\usepackage{amsmath}

\begin{document}

Let us consider two probabilistic events. The first one: Professor M enters the room. The second probabilistic event: Professor M kills a student. Very well then. Which probabilistic events bring more information, more pieces of information?

We have two probabilistic events that happen, you see. The first one: Professor M enters the room. The second one: Professor M kills a student.

Which of these pieces of information would be interesting for mass media?

Student: That Professor M kills the student.

Very well then. This is a question: Why?

Student: Because the second option would imply the first option always.

No, no, we may assume that it is completely different.

Student: But a professor can kill a student without entering the room?

Yes, say on the street, and so on and so forth.

Student: Because this is something unexpected.

Yes. It is the next step, the next logical step.

Student: The second one is less probable.

Yes, it is the first logic, you see. The less probable event, probabilistic event, the more information it gives. It was the first logical step by Shannon.

Information is some function $f$ of inverse probability, okay? Very good.

The second logical step: If you have two probabilistic events which are independent, statistically independent of each other, what can you tell about the joint information about these events?

Quite logically, quite logically, the joint information is the sum of information. If two probabilistic events are independent, yes, then the joint information just has a sum of information, some pieces of information from the first one and from the second one.

The only function, the only function which satisfies such conditions is logarithm. I believe that it is quite clear that logarithm satisfies such conditions. But you may prove the opposite, the converse statement, that the only function which satisfies this condition is logarithm.

So, conclusion: The information is the logarithm of one over the probability, yes? Very good.

Oh, by the way, traditionally one uses natural logarithm. Traditionally one uses natural logarithm. But if you take a binary logarithm with the base 2, you will get your conventional information in bits and bytes, okay? Very good.

That is information. What is entropy?

Let us assume that there are $N$ states of the system. And probability to find itself in the $i$-th state is $p_i$. Probability to find itself in the $i$-th state is $p_i$.

Shannon introduced an extremely interesting measure. Actually, it is an average information about the system in question. The average information about the system in question.

You may consider this concept in two ways.

The first one: You know about the first idea. You know that the system can be in one of the states with probability $p_i$. After that, you observe the system and find it in some particular state, say the first one. It means that you receive some information about the system, yes? If you find the system in this very particular state, then you get this very piece of information about the system.

The question: How many information you have received, yes? It means that you now have this information about the system, yes? It was measured, and you respectively get this number of information. This quantity of information, this value of information. Is it clear?

On the other hand, on the other hand, if you don’t… If you know this probabilistic information about the system and you did not observe the system, you know nothing about the system at hand except these probabilities. It means that system is somewhat mysterious to you. Of course, you know something about it, but it is not complete knowledge, because otherwise if it was complete knowledge, you would say definitely the system is in such and such state. Then your knowledge about system is incomplete, yes?

What is the measure of this uncompleteness? What is the measure of system uncertainty? That’s entropy. Is it clear?

So far, so good? Because it’s rather straightforward material.

You may analyze rigorously this expression, you may analyze rigorously this expression, and find that its minimum is zero. Its minimum is zero. And more to the point, it achieves its minimum when one of—say the first—states, probability for the first state is equal to one, and probability for all other states is equal to zero.

This means that, actually, all terms but the first one would be equal to zero, because $p_i$ if $p_1 = 1$ is equal to zero, and the first term would be equal to zero, because logarithm of one is zero, yes?

What does it mean? It means that if you know… If you knew nothing about the system, its entropy is equal to zero. The measure of its uncertainty is equal to zero.

On the other hand…

\end{document}